\section{Experiment 3: Float versus Double (Hardware FPU vs. Soft-Float)}

\subsection{Aim}
To evaluate the execution speed of single-precision (\texttt{float}) versus double-precision (\texttt{double}) arithmetic, determine the hardware capabilities of the ESP32's floating-point unit (FPU), and uncover the performance hazards of untyped floating-point constants.

\subsection{Background}
The Tensilica Xtensa LX6 architecture implemented in the ESP32 contains a hardware single-precision floating-point coprocessor conforming to ANSI/IEEE Std 754-1985. It includes dedicated 32-bit floating-point registers (f0--f15) and hardware instructions for single-precision addition, subtraction, multiplication, and fused multiply-accumulate. Crucially, the ESP32 lacks a double-precision (64-bit) FPU. Any operation involving a \texttt{double} operand must be synthesized in software using GCC runtime emulation routines (such as \texttt{\_\_adddf3}, \texttt{\_\_muldf3}, and \texttt{\_\_divdf3}). In C and C++, untyped floating-point literals (e.g., \texttt{0.5}) are strictly treated as \texttt{double} by default, whereas literals suffixed with \texttt{f} (e.g., \texttt{0.5f}) are treated as \texttt{float}.

\subsection{Prediction (Recorded Prior to Measurement)}
% TODO(Team): Record your hypothesis here BEFORE running the firmware.
% Expected float vs double cycle counts for add/mul/div/sqrt, and the
% expected penalty of omitting the 'f' suffix on a literal.
\textit{[Pending: team hypothesis not yet recorded.]}

\subsection{Method}
We instantiated volatile \texttt{float} and \texttt{double} variables and benchmarked addition, multiplication, division, and square root across 10 trials of 100 unrolled operations. We specifically evaluated the literal performance penalty by comparing \texttt{f = f * 0.5f} directly against \texttt{f = f * 0.5}.

\subsection{Observations}
% Data source: ../data/exp3_floats.txt (raw Serial Monitor capture from speed_of_numbers.ino, command '3')
\textit{[Pending: no hardware measurements recorded yet. Run Experiment 3 on the ESP32 and populate \texttt{data/exp3\_floats.txt}.]}

\subsection{Analysis}
\textit{[Pending: analysis withheld until measurements are collected.]}

\subsection{What We Learnt}
\textit{[Pending.]}

\subsection{LLM Log}
\begin{llmbox}[LLM Prediction Test: Double Division Latency]
\textbf{Prompt to LLM:} \textit{[Pending: record the verbatim prompt given to the LLM.]} \\
\textbf{LLM Response Summary:} \textit{[Pending.]} \\
\textbf{Board Agreement:} \textit{[Pending.]}
\end{llmbox}
